% Vista pública derivada de 73_rev2_electron_lectores_torres.tex: cuerpo exacto y único retítulo académico del lector temporal de 108 estados. El fragmento del handoff permanece byte-exacto.
\section{Construction complète de l'électron bêta : parcours central, signature, caractère et masse de référence}
\Needspace{7\baselineskip}
\subsection{Parcours central de l'électron et séquence d'enregistrements généalogiques}
L'électron occupe l'unique centre \((5,5)\), avec l'orientation active \(++\) :

\[
 \Gamma_e=(5,5,++),\qquad
 \tau_{12}=+++---+++---.
\]
Ses lecteurs de déplacement sont

\[
 (D_1,D_3,D_4,D_{27},D_{36})=(54,56,68,48,32).
\]
Il ne faut pas confondre :

\begin{enumerate}
\item la signature brute \((24,0,12,0,-12)\) ;
\item la signature paire CT--108 \((-12,-4,12,-6,12)\) ;
\item le vecteur de douze événements ;
\item l'origine affine \(v_e=0\) du caractère relatif.
\end{enumerate}
Il ne faut pas non plus confondre \(\Gamma_e\) avec le mot \(w_e=201101\) du lecteur de la constante d'Euler \(e\).

\Needspace{7\baselineskip}
\subsection{Sélecteur de base de l'électron et transition vers la valeur bêta de référence}
Le sélecteur interne, construit sans utiliser la masse observée, établit

\[
 -22=-\left|((\mathbb Z/9\mathbb Z)\times\mathbb F_3)\setminus\iota(R_\pi)\right|,
 \qquad
 +15=|R_\pi\times\mathbb F_3|.
\]
La lecture des lacunes fournit en outre

\[
 V_e=\begin{pmatrix}21&22\\6&7\end{pmatrix},\qquad
 \det V_e=15.
\]
Le signe négatif de 22 enregistre le déficit/la coupure ; le signe positif de 15, la rétention/l'aire orientée. La matrice est locale à l'électron et ne s'étend pas en sélecteur universel.

L'étape de base est

\[
 e_0=\frac{\sqrt3}{4}
 \left[\exp\!\left(\frac{\varphi}{\pi^2}\right)
 -22\alpha_{\rm HMT}^3\right](1+15\Delta_4),
\]
\[
 e_0=0.510998922488\ldots\ \mathrm{MeV}.
\]
Les généalogies entières sont

\[
 22=\frac{396}{18}=27-5=24-2,
 \qquad
 15=\frac{270}{18}=\binom62=\binom64.
\]
\Needspace{7\baselineskip}
\subsection{Lecteurs bidirectionnels et bêta}
Pour tout parcours,

\[
 K_D(\Gamma)=\left(D_{27}+D_{36},D_1,\frac{D_4-D_3}{2}\right),
\]
\[
 K_\Omega(\Gamma)=(\Omega_{90,120},54,P_6).
\]
Sur les 324 parcours apparaissent 14 profils \(K_D\), 9 profils \(K_\Omega\), 40 paires conjointes et 18 coïncidences. L'unique valeur commune est

\[
 (80,54,6).
\]
Ce 80 procède du lecteur bidirectionnel du parcours électronique, et non de la période 80 du mot d'Euler. Ainsi

\[
 \kappa_e=80-54\alpha_{\rm HMT}+6\Delta_4,
\]
\[
 \boxed{
 \mathfrak e_e^\beta=e_0R_{\rm act}^{\kappa_e}
 =0.5109989506900008086082571648\ldots\ \mathrm{MeV}.}
\]
C'est le résultat électronique de référence. La valeur de base ne le remplace pas. L'égalité avec la masse de l'antiparticule procède de la conjugaison du parcours et de la parité de l'opérateur de masse, et non de l'ajout d'une seconde ligne numérique.

\Needspace{7\baselineskip}
\section{Réalisation physique par couplage spectral}
\Needspace{7\baselineskip}
\subsection{Obstruction à la sélection ponctuelle équivariante d'une multisection non centrale}
Soit \(F\subset\mathcal R_{324}\) une fibre de parcours compatible avec une classe structurelle. Son espace linéaire est

\[
 \mathcal H_F=\bigoplus_{\Gamma\in F}\mathbb C e_\Gamma.
\]
Le caractère central et les lecteurs définissent des opérateurs diagonaux :

\[
 \widehat{\mathcal A}_F^{(0)}e_\Gamma
 =\chi_{\rm full}(\sigma_\Gamma,\eta_\Gamma)e_\Gamma,
\]
\[
 \widehat K_F^r e_\Gamma=\kappa_r(\Gamma)e_\Gamma,
 \qquad r\in\{D,\Omega\}.
\]
La correction two--way est

\[
 \widehat{\mathcal M}_F^{\beta,r}
 =R_{\rm act}^{\widehat K_F^r/2}
  \widehat{\mathcal M}_F^{(0)}
  R_{\rm act}^{\widehat K_F^r/2}.
\]
Lorsque \(\dim\mathcal H_F>1\), ces opérateurs ont généralement plusieurs valeurs propres. En choisir une parce qu'elle est proche d'une masse connue inverserait la causalité. Il n'est pas non plus correct d'exiger qu'un nom physique désigne une cellule unique : l'espèce non centrale réside dans une représentation de la fibre.

\Needspace{7\baselineskip}
\subsection{Foncteur spectral de la réalisation physique : noyau positif de couplage, sous-espace persistant et mesure spectrale}
Une préparation physique \(P\) fournit des données de représentation, et non une masse :

\[
 \mathfrak r(P)=
 (q,\,j,\,\epsilon_{2\pi},\,\epsilon_{4\pi},\,
  \text{statistique},\,\text{génération},\,\text{saveur},\,
  \text{orbite de couleur},\,\text{canal d'interaction}).
\]
Soit \(\mathfrak M\) un couplage compatible avec cette représentation. On définit un noyau positif

\[
 K_{P,\mathfrak M}:F\times F\longrightarrow\mathbb C
\]
qui satisfait :

\begin{enumerate}
\item \(K=K^*\) y \(K\ge0\);
\item \(K(\Gamma,\Gamma')=0\) si les parcours violent la charge, la parité, la couleur, la statistique ou
le canal de \(P\) ;

\item \(K\) commute avec les symétries que le couplage ne brise pas ;
\item \(K(C_M\Gamma,C_M\Gamma')=\overline{K(\Gamma,\Gamma')}\) par conjugaison ;
\item ses entrées sont calculées à partir de l'enregistrement généalogique, de la frontière, de la mémoire, de la retenue et de l'holonomie, sans
grandeurs externes.

\end{enumerate}
L'opérateur d'évolution couplé est

\[
 \mathcal U_{P,\mathfrak M}
 =K_{P,\mathfrak M}^{1/2}\,\mathcal U_{\rm TPK}\,
  K_{P,\mathfrak M}^{1/2}.
\]
Le sous-espace persistant s'obtient par projection spectrale sur les modes qui survivent au retour et à l'élagage :

\[
 \Pi_{P,\mathfrak M}^{\rm surv}
 =\mathbf1_{\mathcal S_{\rm surv}}(\mathcal U_{P,\mathfrak M}).
\]
La réalisation physique correcte est le foncteur spectral

\[
 \boxed{
 \mathcal S_{\rm phys}^{\rm spec}(P,\mathfrak M)
 =\left(
 \mathcal H_{P,\mathfrak M}^{\rm surv},
 \widehat{\mathcal M}_{P,\mathfrak M},
 \mu_{P,\mathfrak M}
 \right),}
\]
où

\[
 \mathcal H_{P,\mathfrak M}^{\rm surv}
 =\Pi_{P,\mathfrak M}^{\rm surv}\mathcal H_F,
\]
\[
 \widehat{\mathcal M}_{P,\mathfrak M}
 =\Pi_{P,\mathfrak M}^{\rm surv}
  \widehat{\mathcal M}_F
  \Pi_{P,\mathfrak M}^{\rm surv},
\]
et \(\mu_{P,\mathfrak M}\) est sa mesure spectrale dans l'état préparé.

Lorsque le couplage est spécifié par une représentation irréductible \(\lambda\) d'un groupe fini \(G_P\), le projecteur correspondant s'obtient sans choisir de parcours distingué :

\[
 P_{P,\lambda}
 =\frac{\dim\lambda}{|G_P|}
  \sum_{g\in G_P}\overline{\chi_\lambda(g)}\,U(g).
\]
Le résultat bidirectionnel est alors la paire comprimée

\[
 \mathbb M(P,\lambda)=
 \left(
 P_{P,\lambda}\widehat{\mathcal M}^{\beta,D}P_{P,\lambda},
 P_{P,\lambda}\widehat{\mathcal M}^{\beta,\Omega}P_{P,\lambda}
 \right).
\]
Si les deux opérateurs commutent avec la représentation irréductible, le lemme de Schur impose leur scalarité sur le bloc. S'ils ne commutent pas, le résultat correct est le spectre du bloc invariant minimal.

\Needspace{7\baselineskip}
\subsection{Naturalité du foncteur spectral sous les entrelaceurs de la dynamique TPK, du couplage et de l'opérateur de masse}
Si \(V:F\to F'\) entrelace la dynamique TPK, le noyau de couplage et les opérateurs de masse, alors

\[
 V\Pi_{P,\mathfrak M}^{\rm surv}
 =\Pi_{P',\mathfrak M'}^{\rm surv}V,
 \qquad
 V\widehat{\mathcal M}_{P,\mathfrak M}
 =\widehat{\mathcal M}_{P',\mathfrak M'}V.
\]
Par le calcul fonctionnel, la mesure spectrale est transportée par \(V\). Le résultat ne dépend ni d'une numérotation des cellules ni du choix d'une base au sein d'une orbite. C'est la condition qui doit remplacer tout sélecteur nominal non équivariant.

\Needspace{7\baselineskip}
\subsection{Conditions d'existence d'une masse scalaire sous le lecteur dimensionnel}
Une masse scalaire est déterminée dans chacun des cas suivants :

\begin{enumerate}
\item le sous-espace persistant est de rang un ;
\item une représentation irréductible fixe un unique projecteur spectral ;
\item le protocole de mesure définit un état positif normalisé
\(\omega_{P,\mathfrak M}\) sur l'algèbre des observables ;

\item une résonance définit un pôle isolé de la résolvante.
\end{enumerate}
Dans le troisième cas,

\[
 m(P,\mathfrak M)
 =\omega_{P,\mathfrak M}
  (\widehat{\mathcal M}_{P,\mathfrak M}).
\]
En l'absence de l'un de ces mécanismes, le résultat scientifique est l'opérateur, son spectre et ses multiplicités. Ce résultat n'est pas incomplet : c'est l'objet déterminé par la dynamique.

\Needspace{7\baselineskip}
\subsection{Construction finie du noyau de masse par filtres de représentation, de persistance et de couplage, suivie d'une normalisation irréductible}
Le noyau s'obtient en quatre étapes finies :

\begin{enumerate}
\item \textbf{Filtre de représentation.} On retient les parcours dont le secteur, la branche de charge,
la couleur, l'holonomie \(2\pi/4\pi\), la parité et la statistique coïncident avec \(P\).

\item \textbf{Filtre de persistance.} On utilise l'ordre partiel formé par le retour de classe,
le retour d'orbite, le retour de cellule, la séparation des préfixes, la dette de retenue et la stabilité décimale. Le résultat peut être une orbite de Pareto ; on n'impose pas l'unicité.

\item \textbf{Filtre de couplage.} On évalue l'incidence de frontière et le canal
d'interaction. Les parcours découplés reçoivent un poids nul.

\item \textbf{Normalisation.} Sur chaque composante irréductible, la trace du
noyau est normalisée à un.

\end{enumerate}
L'atlas orienté fournit toutes les données des deux premières étapes. La troisième exige que chaque réalisation physique déclare son canal ; la quatrième est canonique une fois les composantes irréductibles fixées.

\Needspace{7\baselineskip}
\section{Transducteur généalogique des tours}
\Needspace{7\baselineskip}
\subsection{Paire de mots associés à un même parcours généalogique}
Chaque \(\Gamma\in\mathcal R_{324}\) détermine simultanément un mot temporel

\[
 w_\Gamma\in\mathbb F_3^{108}
\]
et un mot dodécaphasique signé

\[
 \tau_\Gamma\in\{-1,0,+1\}^{12}.
\]
Le premier conserve les déplacements le long de CT--108 ; le second conserve l'orientation accumulée des douze événements. Les deux mots sont liés par le même enregistrement généalogique, mais ne sont pas interchangeables. Le raffinement des tours lit les deux et conserve, en outre, feuillet, retenue, enroulement et mémoire de retour.

\Needspace{7\baselineskip}
\subsection{Correspondance entre hauteur harmonique et longueur temporelle}
Écrivons

\[
 \mathfrak S=30\langle3,4\rangle
 =\{90,120\}\cup\{30r:r\ge6\}.
\]
Pour \(h=30r\), la longueur temporelle associée est

\[
 \ell_h=9r.
\]
En particulier,

\[
 90\longleftrightarrow27,\qquad
 120\longleftrightarrow36,\qquad
 360\longleftrightarrow108.
\]
Cette correspondance n'identifie pas la tour au temps : elle entrelace la hauteur de l'opérateur harmonique avec la longueur du segment de parcours qui l'engendre.

\Needspace{7\baselineskip}
\subsection{Lecteur de déplacement du cycle temporel de 108 états}
Nous définissons l'écart cyclique

\[
 D_\Gamma(n)=
 \sum_{t=0}^{107}
 \mathbf1_{\{w_{\Gamma,t}\ne
 w_{\Gamma,t+n\!\!\!\pmod{108}}\}}.
\]
Le coefficient relatif de la branche temporelle est

\[
 \boxed{\eta_{30r}^{D}(\Gamma)
 =D_\Gamma(9r)-D_{\Gamma_e}(9r).}
\]
La soustraction n'utilise pas une masse comme origine : elle utilise le parcours central comme origine affine du caractère relatif. On a

\[
 \eta_{90}^{D}+\eta_{120}^{D}
 =[D_\Gamma(27)+D_\Gamma(36)]-80,
\]
de sorte que la première couche prolonge exactement la première composante du lecteur \(K_D\).

\Needspace{7\baselineskip}
\subsection{Lecteur d'accumulation dodécaphasique}
Pour \(r\ge1\), soit

\[
 A_\Gamma(r)=
 \sum_{j=0}^{11}
 \left(
 \sum_{u=0}^{r-1}
 \tau_{\Gamma,j+u\!\!\!\pmod{12}}
 \right)^2.
\]
Le coefficient relatif de la branche orientée est

\[
 \boxed{\eta_{30r}^{\Omega}(\Gamma)
 =A_\Gamma(r)-A_{\Gamma_e}(r).}
\]
Alors

\[
 4\eta_{90}^{\Omega}-3\eta_{120}^{\Omega}
 =\Omega(\tau_\Gamma)-80,
\]
de sorte que cette couche prolonge la première composante de \(K_\Omega\). Les deux familles forment le relèvement bidirectionnel

\[
 \mathcal F_{\rm tower}^{2w}:
 \mathcal R_{324}\longrightarrow
 \mathcal E_{\mathfrak S}\times\mathcal E_{\mathfrak S},
 \qquad
 \Gamma\longmapsto
 \bigl(\eta^D(\Gamma),\eta^\Omega(\Gamma)\bigr).
\]
\Needspace{7\baselineskip}
\subsection{Détermination finie de la tour infinie}
La périodicité du mot temporel donne

\[
 D_\Gamma(9(r+12))=D_\Gamma(9r).
\]
Si

\[
 T_\Gamma=\sum_{j=0}^{11}\tau_{\Gamma,j},
 \qquad r=12q+s,\quad 0\le s<12,
\]
alors

\[
 A_\Gamma(12q+s)
 =A_\Gamma(s)+(12q^2+2qs)T_\Gamma^2.
\]
Ainsi, toute la tour est déterminée par douze valeurs de \(D_\Gamma\), douze valeurs de \(A_\Gamma\) et l'entier \(T_\Gamma^2\). L'infinité des hauteurs n'exige pas une infinité de données indépendantes.

La même formule résout le diamant généalogique de hauteur 360. Si

\[
 h(a,b)=90a+120b,
\]
alors

\[
 (a+4,b)\sim(a,b+3)
\]
puisque \(3(a+4)+4b=3a+4(b+3)\). La généalogie \((a,b)\) est conservée jusqu'à l'application de la mémoire et de la retenue ; ce n'est qu'ensuite qu'elle descend à la hauteur commune.

\Needspace{7\baselineskip}
\subsection{Convergence absolue du caractère raffiné sur les tours \(90/120\)}
On a

\[
 |\eta_h^D|\le108,\qquad
 |\eta_{30r}^{\Omega}|\le24r^2.
\]
Comme \(0<q_+<q_-<1\),

\[
 |s_{30r}|\le2q_-^{30r}.
\]
On en déduit

\[
 \sum_{h\in\mathfrak S}|\eta_h^D s_h|<\infty,
 \qquad
 \sum_{h\in\mathfrak S}|\eta_h^\Omega s_h|<\infty.
\]
Le caractère raffiné est bien défini sans tronquer la tour.

\Needspace{7\baselineskip}
\subsection{Mémoire, retenue et cocycle}
Le coefficient brut est enrichi par

\[
 \mathbb F_{\rm tower}(\Gamma)_h
 =\bigl(\eta_h^D,\eta_h^\Omega,\rho_h,
 \mathbf w_h,c_h,N_h^{\rm surv}\bigr),
\]
où la signature tronquée satisfait

\[
 \sigma_{\Gamma|\ell_h}
 =r(\rho_h)+\mathcal D\mathbf w_h,
 \qquad
 \mathcal D=\operatorname{diag}(2,6,2,2,4),
\]
et la retenue de deux résidus est

\[
 c(\rho_1,\rho_2)=
 \mathcal D^{-1}
 \bigl[r(\rho_1)+r(\rho_2)-r(\rho_1+\rho_2)\bigr].
\]
Le défaut de localité d'une composition,

\[
 \delta_h(\Gamma,\Lambda)=
 F_h(\Gamma\star\Lambda)-F_h(\Gamma)-F_h(\Lambda),
\]
obéit à l'identité de cocycle

\[
 \delta_h(\Gamma,\Lambda)
 +\delta_h(\Gamma\star\Lambda,\Xi)
 =\delta_h(\Lambda,\Xi)
 +\delta_h(\Gamma,\Lambda\star\Xi).
\]
On préserve ainsi la raison pour laquelle deux parcours de même hauteur visible peuvent transporter des histoires distinctes.

\Needspace{7\baselineskip}
\subsection{Sélection de la lecture par le couplage}
Les deux branches ne sont pas moyennées par commodité. Le couplage physique \(a\) doit déterminer une transformation naturelle

\[
 \Pi_a:
 \mathcal E_{\mathfrak S}\times\mathcal E_{\mathfrak S}
 \longrightarrow\mathcal E_{\mathfrak S}
\]
qui commute avec les troncatures et les symétries, respecte le cocycle de retenue, conserve l'origine électronique et soit compatible avec \(X_{90}^{4}=Y_{120}^{3}\). Si le couplage échange les deux lectures, l'unique clôture continue, multiplicative, symétrique et normalisée sur les scalaires positifs est la moyenne géométrique ; en termes d'exposants,

\[
 L_{\rm sym}=\frac{L_D+L_\Omega}{2}.
\]
En l'absence de cette symétrie, le résultat est la mesure spectrale conjointe de \((L_D,L_\Omega)\), et non une moyenne arbitraire.

\Needspace{7\baselineskip}
\subsection{Raffinement coinductif par orbites primitives}
Le relèvement fini précédent admet un prolongement. Soit \(\mathscr X_\Gamma\) le graphe des extensions admissibles du parcours, \(U_\Gamma\) son opérateur de transfert et \(R_\Gamma\) l'involution d'orientation. Les traces

\[
 a_n(\Gamma)=\operatorname{Tr}(R_\Gamma U_\Gamma^n)
\]
et l'inversion de Möbius

\[
 p_n(\Gamma)=\frac1n\sum_{d\mid n}\mu(d)a_{n/d}(\Gamma)
\]
séparent les orbites primitives. La différence

\[
 p_h(\Gamma)-p_h(\Gamma_e)
\]
est un raffinement ultérieur du coefficient de hauteur \(h\), et non un remplacement des deux lectures finies.

\Needspace{7\baselineskip}
\subsection{Formulation historique au moyen d'un automate d'extension}
La signature \(\mathbb Z^5\) n'épuise pas l'histoire d'un parcours. Pour chaque \(\Gamma\in\mathcal R_{324}\), on conserve :

\begin{itemize}
\item les 108 transitions de l'enregistrement généalogique ;
\item les vingt blocs et leurs préfixes ;
\item les mots \(w_6\) ;
\item les neuf comparaisons \(K\leftrightarrow K+9\) ;
\item le feuillet, la retenue et l'enroulement ;
\item les premiers retours de classe, d'orbite, de cellule et d'état cinématique ;
\item l'ambiguïté accumulée et les triades décimales stabilisées ;
\item les profils \(K_D\) et \(K_\Omega\).
\end{itemize}
Ces données permettent de construire un transducteur de tours sans consulter le résidu d'une masse.

\Needspace{7\baselineskip}
\subsection{Automate local d’extension}
Soit \(\mathscr X_\Gamma\) l’ensemble fini des états locaux compatibles avec la route. Un état conserve

\[
 x=(t,g,B,p,\varepsilon,c,w),
\]
où \(t\) est le temps discret, \(g\) la phase nonadique, \(B\) le bloc, \(p\) le préfixe, \(\varepsilon\) le feuillet, \(c\) la retenue et \(w\) l'enroulement. Il existe une arête \(x\to y\) lorsque \(y\) est une extension admissible de \(x\), respecte l'enregistrement généalogique et satisfait l'élagage \(G_9\). Le tableau des 36.864 compositions de retenue fournit la loi locale d'addition ; les tableaux de retour fixent la mémoire.

Soit \(U_\Gamma\) l’opérateur de transfert normalisé de ce graphe et \(R_\Gamma\) l’involution d’orientation. On exige

\[
 U_{C_M\Gamma}=J U_\Gamma J^{-1},\qquad
 R_{C_M\Gamma}=-J R_\Gamma J^{-1}.
\]
\Needspace{7\baselineskip}
\subsection{Dénombrements orientés de cellules primitives et construction de la signature électronique}
Pour \(n\ge1\), la trace orientée

\[
 a_n(\Gamma)=\operatorname{Tr}(R_\Gamma U_\Gamma^n)
\]
dénombre les retours fermés avec signe. L’inversion de Möbius sépare les orbites primitives :

\[
 p_n(\Gamma)
 =\frac1n\sum_{d\mid n}\mu(d)a_{n/d}(\Gamma).
\]
On prend l’électron comme origine affine des corrections relatives et on définit

\[
 \boxed{
 \eta_h(\Gamma)=p_h(\Gamma)-p_h(\Gamma_e),
 \qquad h\in\mathfrak S.}
\]
On obtient ainsi l’application

\[
 \mathcal F_{\rm tower}:
 \Gamma^{\rm enr}\longmapsto
 \eta(\Gamma)=(\eta_h(\Gamma))_{h\in\mathfrak S}.
\]
Il n’y a pas de paramètres par espèce. Le même automate, la même involution et la même inversion de Möbius agissent sur les 324 routes.

\Needspace{7\baselineskip}
\subsection{Convergence du raffinement primitif}
Soit \(N_\Gamma=\dim\mathscr X_\Gamma\). Si \(U_\Gamma\) est une contraction,

\[
 |a_n(\Gamma)|\le N_\Gamma.
\]
Le nombre de diviseurs de \(n\) croît de manière sous-linéaire, de sorte que

\[
 |p_n(\Gamma)|
 \le\frac{N_\Gamma\tau(n)}n.
\]
Comme \(0<q_+<q_-<1\),

\[
 |s_h|\le q_-^h+q_+^h\le2q_-^h.
\]
Par comparaison,

\[
 \sum_{h\in\mathfrak S}|\eta_h(\Gamma)s_h|<\infty
\]
uniformément sur toute collection finie de fibres. Le caractère raffiné est donc bien défini.

\Needspace{7\baselineskip}
\subsection{Identités de conservation du parcours électronique et unicité du point fixe central}
La construction doit satisfaire simultanément :

\begin{enumerate}
\item invariance par renumérotation des états ;
\item covariance sous conjugaison de feuillet ;
\item compatibilité avec la somme directe de composantes ;
\item égalité lors du recalcul à partir du registre généalogique et du graphe certifié ;
\item conservation de l’origine \(\eta(\Gamma_e)=0\) ;
\item insensibilité absolue à une permutation des colonnes externes de masse ;
\item détermination des queues au moyen d’une borne antérieure à la confrontation.
\end{enumerate}
La représentation exacte d’un résiduel donné démontre la complétude du codomaine, mais ne remplace pas ces sept preuves.

\Needspace{7\baselineskip}
\section{Horloge de route et normalisation absolue}
\Needspace{7\baselineskip}
\subsection{Les retours entiers ne suffisent pas à distinguer les espèces}
Dans l’atlas de 324 routes, le premier retour de cellule est 27, le retour cinématique au sein de CT--108 est 54 et le retour cinématique étendu est 216. Étant communs, ces entiers fixent l’architecture temporelle globale, mais ne peuvent pas expliquer à eux seuls les différences de masse entre espèces.

\Needspace{7\baselineskip}
\subsection{Holonomie temporelle du parcours électronique et retour avec mémoire}
L’information spécifique réside dans l’holonomie de l’opérateur d’extension. Soit \(\lambda_j(\Gamma)\) une valeur propre non triviale de \(U_\Gamma\) sur le sous-espace persistant. Nous écrivons

\[
 \lambda_j(\Gamma)=r_j(\Gamma)e^{i\theta_j(\Gamma)}.
\]
La fréquence interne associée est

\[
 \omega_j(\Gamma)
 =\frac{\theta_j(\Gamma)+2\pi w_j(\Gamma)}{108t_0},
\]
où l’entier \(w_j\) est fixé par l’enroulement enregistré, non par un choix de branche ultérieur. Pour un mode stable \(r_j=1\) ; pour une résonance \(r_j<1\) et son atténuation apporte une largeur.

Le quotient

\[
 \rho_{\omega,j}(\Gamma)
 =\frac{\omega_j(\Gamma)}{\omega_0}
\]
est adimensionnel et peut modifier les rapports de masse sans altérer la décade commune \(-34\). L’énergie et la masse du mode sont

\[
 E_j(\Gamma)=\hbar_{\rm ret}\omega_j(\Gamma)
 \chi_{\rm full}(\Gamma)R_{\rm act}^{\kappa_j(\Gamma)},
\]
\[
 m_j(\Gamma)=E_j(\Gamma)c_{\rm int}^{-2}.
\]
Cette équation réunit action, horloge, caractère et lecteur sans convertir l'un en l'autre. Son application numérique exige de déterminer \(U_\Gamma\), de fixer l'enroulement et de déclarer la carte dimensionnelle.

\Needspace{7\baselineskip}

